% ID: SEL-D0016-BD-E08
% ESTADO: cotejada. Segunda pasada visual completa y borrador compilado revisados.
% Fuente: Selectividad/matematicas_aplicadas/matematicas_aplicadas_P_D.pdf
% SHA-256: 11865b83e509d9cc2723c752f7a81f714378e1b651f4a4d010dab4a4576b74c2
% Página: 2; bloque D; ejercicio original 8.

El tiempo de espera para ser atendido en un servicio hospitalario es una
variable aleatoria que sigue una distribución Normal con desviación típica
de 2 meses. Tomada una muestra al azar de 9 pacientes que han utilizado ese
servicio, se han registrado los siguientes tiempos de espera en meses:
\[
8.5\quad 3.7\quad 4.3\quad 3.6\quad 5.6\quad 4.8\quad 1.0\quad
1.4\quad 6.0
\]
\begin{enumerate}[label=\alph*)]
  \item \textbf{(1.5 puntos)} Determine un intervalo de confianza al 95\%
  para el tiempo de espera medio poblacional.
  \item \textbf{(1 punto)} Con un nivel de confianza del 97\%,
  \textquestiondown{}qué tamaño muestral mínimo se ha de tomar para que
  el error máximo cometido en la estimación del tiempo de espera medio
  poblacional no exceda de un mes?
\end{enumerate}
